Enfin, pour la vitesse tout au long de la glissade, les équations de la
conservation de l'énergie sont utilisées, ce qui nous permet de décrire notre
système comme suit:

\begin{align}
\underbrace{m g y_0}_{\textrm{énergie totale}} =
  \underbrace{m g F(x)}_{\textrm{énergie potentielle à } x} +
  \underbrace{\frac{1}{2} m [v(x)]^2}_{\textrm{énergie cinétique à } x} +
  \underbrace{\int_0^x F_f(x) \, ds}_{\textrm{travail cumulatif de la friction à } x}
\end{align}

La trajectoire obtenue permet d'établir un angle de pente en radians selon $x$:

\begin{equation}
\theta(x) = \arctan(F'(x))
\end{equation}

La force de frottement le long de la trajectoire est donnée par $F_f(x) = \mu_f
m g \cos\theta(x)$, et le travail de frottement s'écrit correctement en fonction
de l'arc le long de la trajectoire:

\begin{equation}
\int_0^x F_f(x) \, ds = \mu_f m g \int_0^x \cos\theta(x) \sqrt{1 + [F'(x)]^2} \, dx
\end{equation}

Sous l’hypothèse que la force normale est exactement $mg\cos\theta(x)$ et que le
frottement agit tangentiellement à la trajectoire, le facteur $\sqrt{1 +
[F'(x)]^2}$ du déplacement le long de la pente se simplifie avec
$\cos\theta(x)$, ce qui conduit à:

\begin{equation}
\int_0^x F_f(x) \, ds = \mu_f m g \, x
\end{equation}

c’est-à-dire que le travail de frottement dépend uniquement de la distance
horizontale parcourue.

La vitesse peut alors être isolée:
\begin{equation}
v(x) = \sqrt{2 g \left( y_0 - F(x) - \mu_f \, x \right)}
\end{equation}

En utilisant MATLAB pour recalculer l'intégrale à chaque valeur de $x$ le long
de la trajectoire, il est possible de tracer plusieurs courbes en posant des
$\mu_f$ répartis dans la plage que la valve permet. Lors de la génération des
courbes, il devient apparent que les $\mu_f$ aux extrêmes donnent des vitesses
finales hors des limites acceptables (20-25 km/h). La figure \ref{fig:01-speeds}
ne montre que des valeurs de $\mu_f$ permettant des vitesses finales proches de
la plage d'intérêt.

Un coefficient de friction de $0.62$ est choisi puisqu'il donne, parmi les
trajectoires calculées, une vitesse finale proche du centre de la plage (22.5
km/h), ce qui permet d'accommoder les erreurs provenant de l'approximation.

Ayant choisi un $\mu_f$, il reste à déterminer quelle ouverture de valve permet
de l'atteindre en isolant $x$ dans l'équation \ref{eq:valve-quad} et en
choisissant la racine située entre 0 et 100 (plage physique de la valve).
L'ouverture trouvée est de 33.19. À cette ouverture, une variation de $\pm$RMS
permet un coefficient de friction entre 0.6020 et 0.6380, comme représenté à la
figure \ref{fig:01-speed-comparison}.
